\section{训练阶段与数据对象}

数据工作通常不止一阶段。比较清晰的划分是：

\begin{table}[H]
\centering
\begin{tabular}{p{0.18\textwidth}p{0.28\textwidth}p{0.45\textwidth}}
\toprule
\textbf{阶段} & \textbf{目标} & \textbf{典型数据} \\
\midrule
Pre-training & 学到语言分布与通用知识 & 大规模原始网页、书籍、百科、代码、论文 \\
Mid-training & 纠偏和补强某些能力 & 更高质量网页、数学、代码、长上下文、合成任务 \\
Post-training & 对齐成可对话、可指令跟随的模型 & 指令数据、聊天数据、偏好数据、RL 数据 \\
\bottomrule
\end{tabular}
\caption{语言模型训练常见分阶段}
\end{table}

\begin{knowledgebox}{Base model 与 Instruct model}
\textbf{Base model} 通常指 pre-training 加上 mid-training 后的检查点。\\
\textbf{Instruct/chat model} 通常指 post-training 后的模型，也就是能更自然地跟人对话的版本。
\end{knowledgebox}

讲者强调，阶段边界在现实里并不总是严格清晰。很多模型会插入更多子阶段，但方向都一样：先用大规模低质量数据打底，再用更小规模高质量数据修正和对齐。以 AI2 的 OLMo 为例，其三阶段分别使用了不同的数据集（预训练用 Dolma，mid-training 用 Dolmino Mix，后训练用 T\"ulu 数据），每阶段的数据规模和质量要求差异显著。

\begin{table}[H]
\centering
\begin{tabular}{p{0.16\textwidth}p{0.23\textwidth}p{0.23\textwidth}p{0.27\textwidth}}
\toprule
\textbf{阶段} & \textbf{主要信号} & \textbf{常见风险} & \textbf{工程重点} \\
\midrule
Pre-training & 语言续写、事实共现、风格吸收 & 噪声大、重复多、合规复杂 & 先把规模做起来，再谈更精细的能力 \\
Mid-training & 数学、代码、长文档、网页精选 & 过拟合少数任务、分布偏移 & 用更高质量数据把能力补齐 \\
Post-training & 指令响应、对话、偏好 & 幻觉、拒答不稳、风格失控 & 把模型变成可用的助手，而不是只会续写的文本引擎 \\
\bottomrule
\end{tabular}
\caption{三阶段训练的信号、风险与目标并不相同}
\end{table}

很多人会把“把数据喂得更干净”理解成同一件事，但实际上每个阶段的目标函数都不一样。pre-training 更像让模型见识世界的分布，mid-training 更像补课和纠偏，post-training 则是在控制输出方式和交互方式。若把三者混成一个大混合桶，最后通常只会得到一个既不稳定、也不容易解释的模型。

从数据规模看，三个阶段的量级差异极大：

\begin{table}[H]
\centering
\begin{tabular}{p{0.17\textwidth}p{0.22\textwidth}p{0.22\textwidth}p{0.27\textwidth}}
\toprule
\textbf{阶段} & \textbf{典型 token 量级} & \textbf{数据来源} & \textbf{质量要求} \\
\midrule
Pre-training & 1T -- 36T tokens & 网页、书籍、代码、百科 & 覆盖面优先，噪声可接受 \\
Mid-training & 100B -- 1T tokens & 精选网页、数学、代码 & 质量明显高于预训练 \\
Post-training & 10K -- 1M 样本 & 指令、对话、偏好标注 & 极高质量，逐条审核 \\
\bottomrule
\end{tabular}
\caption{三阶段训练的数据规模差异可达 3--4 个数量级}
\end{table}

